% Generated draft: requires mathematical and source review.
\documentclass[11pt]{article}
\usepackage[margin=25mm]{geometry}
\usepackage{fontspec}
\setmainfont{Latin Modern Roman}
\usepackage{amsmath,amssymb,mathtools,unicode-math}
\setmathfont{Latin Modern Math}
\usepackage{enumitem,needspace,xurl,hyperref}
\setlength{\parindent}{0pt}
\setlength{\parskip}{5pt}
\setlength{\emergencystretch}{3em}
\newcommand{\N}{\mathbb{N}}
\newcommand{\Z}{\mathbb{Z}}
\newcommand{\Q}{\mathbb{Q}}
\newcommand{\R}{\mathbb{R}}
\newcommand{\C}{\mathbb{C}}
\newcommand{\F}{\mathbb{F}}
\newcommand{\A}{\mathbb{A}}
\newcommand{\PP}{\mathbb{P}}
\newcommand{\E}{\mathbb{E}}
\newcommand{\eps}{\varepsilon}
\newcommand{\dd}{\,\mathrm{d}}
\begin{document}
\section*{UnsolvedMath ID: 2307002}
\textbf{Source problem code:} catalogue entry
\subsection{Definitions and mathematical statement}
The source lets $D$ be simply connected and writes a Faber expansion
\[
 f(z)=\sum_{n=0}^{\infty}a_n(f)P_n(z).
\]
A standard Faber construction, when applicable to a bounded domain, uses a normalized conformal map $\Phi$ from the exterior of $\overline D$ to the exterior unit disc, with $\Phi(\infty)=\infty$ and positive leading coefficient, and defines $P_n$ as the polynomial part of $\Phi(z)^n$. The source does not supply the geometric assumptions or normalization needed for its expansion; the displayed contextual construction does not extend that claim to every simply connected domain.

\textbf{Variability problem (W. C. Royster).} For a fixed index $n$, determine the set
\[
 V_n(D)=\{a_n(f):f\text{ analytic in }D,\ \operatorname{Re}f(z)>0\text{ for all }z\in D\},
\]
with the intended Faber basis and any omitted normalization recovered from the source. The statement reports that $V_n(D)$ is a circle when $D$ is a circle, and an ellipse when $D$ is an ellipse, referring to Royster.

As printed, there is no normalization of $f$ to fix scale. Positive real part is preserved by multiplication by any positive constant, so a nonzero attainable coefficient can be scaled arbitrarily. For the unit disc with the usual basis $P_n(z)=z^n$ and $n\geq1$, every $b\in\C$ occurs in $f(z)=c+bz^n$ with $c>|b|$. Consequently the unnormalized reading is incompatible with a bounded circular variability region. A condition such as a prescribed value or real part at a basepoint may be missing, but none is selected here. The index range and whether ``circle/ellipse'' refers to the boundary or filled region are likewise not stated.

The supplied summary records OPEN-TRIAGE, says the problem was open in the 2018 edition, and reports no definitive resolution located. These are status claims about the intended normalized problem, whose full formulation is absent.
\subsection{Short English statement}
Determine the possible Faber coefficients of analytic functions with positive real part on a simply connected domain. The reported circular and elliptical cases indicate missing normalization: as printed, positive scaling makes the coefficient regions unbounded. The basis assumptions and coefficient index also need clarification.
\subsection{Sources}
\begin{itemize}
\item Local UnsolvedMath registry, permanent ID 2307002, source code catalogue entry, ``Research Problems in Function Theory --- Problem 7.2''; attribution: W. C. Royster. Supplied status open, classification OPEN-TRIAGE, updated 2026-08-06.
\item Hayman, \emph{Research Problems in Function Theory} (2018), Problem 7.2, \url{https://arxiv.org/abs/1809.07200}. The record reports source-TeX extraction and access on 2026-07-29; the 2018 update was retained for status triage. Current status and rights require review; the record describes publicly available arXiv source TeX.
\item The statement refers to Royster for the circular and elliptical cases, without a publication title, date, or theorem reference.
\end{itemize}
\end{document}
